\textbf{Start ablation.} Table~\ref{tab:start} gives the two starts separately. Neither dominates. From start S, Jastrow stays far above its start-B result for $\Gamma\le0.75$ (SR at 10/0.5: $3.0\times10^{-3}$ against $7.4\times10^{-5}$) and above the selected MF error in all four $\Gamma\le0.5$ tasks, although the Jastrow family contains the product state; MF-SR at 12/0.5 has a median of $2.3\times10^{-1}$, a domain-wall state, against $1.0\times10^{-3}$ from start B; and all 28 diverged runs are RBM-SR runs from start S. From start B, Jastrow is worse at $\Gamma{=}1$ (SR at 10/1: $4.1\times10^{-3}$ against $1.5\times10^{-3}$) and the RBMs with SR are worse at 10/0.75 ($\alpha{=}2$: $7.6\times10^{-4}$ against $8.1\times10^{-5}$). With start S alone, as registered, H1 also holds on medians, but with margins over Jastrow at small fields that are artefacts of the start. We did not isolate why start B has no divergences.

\begin{table*}[t]\centering\caption{Start ablation: median $\epsilon_E$ over the non-diverged seeds of each start (S: symmetric, B: symmetry-broken; superscript: diverged seeds of 5); columns are $N/\Gamma$. SR systems only, for space.}\label{tab:start}\footnotesize\setlength{\tabcolsep}{3pt}\begin{tabular}{llrrrrrrrrrrrrr}\toprule
System & start & 10/0.25 & 10/0.5 & 10/0.75 & 10/1 & 10/1.25 & 10/1.5 & 10/2 & 8/0.5 & 8/1 & 8/1.5 & 12/0.5 & 12/1 & 12/1.5 \\\midrule
MF SR & S & 6.2e-5 & 1.5e-3 & 6.2e-3 & 2.3e-2 & 5.0e-2 & 6.8e-2 & 6.5e-2 & 1.1e-3 & 2.6e-2 & 6.7e-2 & 2.3e-1 & 2.2e-2 & 6.8e-2 \\
 & B & 6.2e-5 & 1.3e-3 & 6.2e-3 & 2.3e-2 & 5.0e-2 & 6.8e-2 & 6.3e-2 & 1.1e-3 & 2.6e-2 & 6.8e-2 & 1.0e-3 & 2.2e-2 & 6.9e-2 \\
Jas SR & S & 4.0e-4 & 3.0e-3 & 5.3e-3 & 1.5e-3 & 1.9e-4 & 3.2e-5 & 2.6e-6 & 3.3e-3 & 1.2e-3 & 3.7e-5 & 2.8e-3 & 1.7e-3 & 3.0e-5 \\
 & B & 5.6e-5 & 7.4e-5 & 8.4e-4 & 4.1e-3 & 1.9e-4 & 3.3e-5 & 2.6e-6 & 1.4e-4 & 1.6e-3 & 3.7e-5 & 6.0e-5 & 4.3e-3 & 3.1e-5 \\
RBM$\alpha{=}$1 SR & S & 6.1e-5 & 8.7e-5 & 1.2e-4 & 6.3e-5 & 1.3e-5 & 5.9e-6 & 7.2e-7 & 6.2e-5 & 3.3e-5 & 3.2e-6 & 8.8e-5 & 9.4e-5 & 6.4e-6 \\
 & B & 7.3e-5 & 7.7e-5 & 8.0e-4 & 6.7e-5 & 2.0e-5 & 4.2e-6 & 7.2e-7 & 1.2e-4 & 3.7e-5 & 4.0e-6 & 6.0e-5 & 1.4e-4 & 5.9e-6 \\
RBM$\alpha{=}$2 SR & S & 5.2e-5 & 5.2e-5 & 8.1e-5$^{1}$ & 4.0e-5$^{1}$ & 5.9e-6$^{1}$ & 2.2e-6$^{2}$ & 1.4e-6$^{2}$ & 5.8e-5 & 1.7e-5 & 9.6e-7 & 6.1e-5 & 8.0e-5$^{2}$ & 3.0e-6$^{3}$ \\
 & B & 6.7e-5 & 4.4e-5 & 7.6e-4 & 4.5e-5 & 6.1e-6 & 2.0e-6 & 1.7e-6 & 1.2e-4 & 2.7e-5 & 9.9e-7 & 3.7e-5 & 8.3e-5 & 3.2e-6 \\
RBM$\alpha{=}$4 SR & S & 2.0e-5 & 6.1e-5 & 9.7e-5 & 2.3e-5 & 5.4e-6$^{2}$ & 2.5e-6$^{3}$ & 2.1e-6$^{4}$ & 5.2e-5 & 7.2e-6$^{1}$ & 1.1e-6$^{2}$ & 5.6e-5$^{1}$ & 4.5e-5$^{1}$ & 6.1e-6$^{2}$ \\
 & B & 3.7e-5 & 4.1e-5 & 7.7e-4 & 2.7e-5 & 5.7e-6 & 2.3e-6 & 2.4e-6 & 1.2e-4 & 1.5e-5 & 1.1e-6 & 4.7e-5 & 6.5e-5 & 5.8e-6 \\
\bottomrule\end{tabular}
\end{table*}

\textbf{Sampling-free reference.} MF and Jastrow converge within the L-BFGS iteration cap; all RBM runs stop at the cap (Table~\ref{tab:check}), so the RBM row of Table~\ref{tab:energy} is an upper bound that is not tight: at 8/0.5 it is above the VMC median. In the other twelve tasks the RBM-SR median is above the L-BFGS value, by a factor of about 3 at 10/0.5 and about 400 at 10/0.25, so the RBM numbers in this paper measure the optimiser and the sample budget more than the ansatz.

\textbf{SR shift and sample count.} Two sweeps use RBM $\alpha{=}2$ SR at $\eta{=}0.1$ from start S, at 10/1, seeds 0--4; their base setting ($\epsilon{=}10^{-2}$, $M{=}256$) is the start-S main configuration. Table~\ref{tab:sweeps} shows SR is unusable at this rate for shifts $10^{-4}$ and $10^{-3}$ (all five seeds diverged), mostly stable at $10^{-2}$ (1 of 5 diverged), and stable but less accurate at $10^{-1}$ and $1$. We did not retune $\eta$ per shift, so the sweep ties the start-S divergences to the shift--rate combination but does not show that a smaller rate at a small shift would fail. The median $\epsilon_E$ falls at every step from $M{=}32$ ($5.6\times10^{-4}$) to $M{=}512$ ($2.1\times10^{-5}$): the error of the trained RBM is tied to the sample budget, and differences between systems at the $10^{-5}$ level should be read relative to it.

\begin{table}[t]\centering\caption{Sweeps of the SR shift $\epsilon$ (at $M{=}256$) and of the samples per iteration $M$ (at $\epsilon{=}10^{-2}$), RBM $\alpha{=}2$ SR at 10/1 from start S: medians over non-diverged seeds; diverged seeds.}\label{tab:sweeps}\footnotesize\begin{tabular}{lrrrr}\toprule
setting & $\epsilon_E$ & $\epsilon_{M_z}$ & $1{-}F$ & div. \\\midrule
$\epsilon$$\,{=}\,$0.0001 & -- & -- & -- & 5/5 \\
$\epsilon$$\,{=}\,$0.001 & -- & -- & -- & 5/5 \\
$\epsilon$$\,{=}\,$0.01 & 4.0e-5 & 1.2e-3 & 9.1e-5 & 1/5 \\
$\epsilon$$\,{=}\,$0.1 & 7.0e-5 & 1.2e-3 & 1.4e-4 & 0/5 \\
$\epsilon$$\,{=}\,$1 & 3.8e-4 & 1.6e-3 & 7.4e-4 & 0/5 \\
\midrule
$M$$\,{=}\,$32 & 5.6e-4 & 2.5e-3 & 1.6e-3 & 1/5 \\
$M$$\,{=}\,$64 & 2.6e-4 & 4.8e-3 & 7.4e-4 & 0/5 \\
$M$$\,{=}\,$128 & 1.0e-4 & 3.0e-3 & 3.3e-4 & 2/5 \\
$M$$\,{=}\,$256 & 4.0e-5 & 1.2e-3 & 9.1e-5 & 1/5 \\
$M$$\,{=}\,$512 & 2.1e-5 & 3.9e-4 & 3.7e-5 & 0/5 \\
\bottomrule\end{tabular}
\end{table}

\textbf{Training curves.} Fig.~\ref{fig:curves} shows the Monte Carlo energy error during training at 10/1 from start S (median of three seeds). With SR the RBM curves flatten early, whereas RBM-SGD is still decreasing at iteration 300.

\begin{figure}[t]\centering\includegraphics[width=0.66\columnwidth]{curves.pdf}\caption{MC energy error during training (10/1, start S).}\label{fig:curves}\end{figure}

\textbf{Checks and cost.} In Table~\ref{tab:check}, no finite run violates the variational bound, no selected state has $\epsilon_E>0.1$ (13 start-S runs do), and the median gap between MC and exact energy is $2.4\times10^{-5}$ for RBM $\alpha{=}2$ SR.

\begin{table}[t]\centering\caption{Divergence counts, consistency checks, and run counts and times from the registry.}\label{tab:check}\footnotesize\setlength{\tabcolsep}{3pt}\begin{tabular}{lr}\toprule
Quantity & value \\\midrule
diverged of 650 main runs, v1 / v2 (start S) & 31 / 28 \\
diverged of 650, final: start S / B / both & 28 / 0 / 0 \\
of these, RBM$\alpha{=}$2 / 4 SR, start S (of 65) & 12 / 16 \\
$\epsilon_E>0.1$: start S / B / selected & 13 / 0 / 0 \\
min $\epsilon_E$ over the finite main runs & 5.6e-7 \\
median $|\epsilon_E^{MC}-\epsilon_E|$, RBM$\alpha{=}$2 SR & 2.4e-5 \\
tuning runs per SGD / SR system & 36 / 27 \\
L-BFGS runs / stopped at iteration cap & 78 / 26 \\
runs: start S / B / L-BFGS & \rhval{count/main/runs} / \rhval{count/main_b/runs} / \rhval{count/exact_opt/runs} \\
runs: tuning / shift / samples / curves & \rhval{count/tune_lr3/runs} / \rhval{count/sweep_shift/runs} / \rhval{count/sweep_samples/runs} / \rhval{count/curves/runs} \\
summed run time (min): reported / all & 47 / 82 \\
wall-clock with a run executing (min) & 47 \\
longest single run (s) & 115 \\
\bottomrule\end{tabular}
\end{table}
